% !TeX root = ../../apuntes.tex
\section{Introducción a la Flexión}

\subsection{Recordatorio de Esfuerzos Internos en el Prisma Mecánico}

Al practicar un corte transversal en un prisma mecánico en equilibrio y reducir el sistema de fuerzas internas al centro de gravedad $G$ de la sección recta $\Omega$ (de normal exterior $\vec{i}$ según la directriz longitudinal $x$), obtenemos el vector resultante de fuerzas $\vec{R}$ y el momento resultante $\vec{M}$:
\begin{align}
    \vec{R} & = N\,\hat{i} + Q_y\,\hat{j} + Q_z\,\hat{k} \\
    \vec{M} & = T\,\hat{i} + M_y\,\hat{j} + M_z\,\hat{k}
\end{align}

\begin{figure}[H]
    \centering
    \adjustbox{max width=0.92\textwidth}{%
        \begin{tikzpicture}[>=Latex, scale=1.0]
            % 1. Prisma cilíndrico con reducción de esfuerzos en G
            \begin{scope}[shift={(0, 2.2)}]
                \node[above, font=\small\bfseries, blue!85!black] at (2.2, 2.75) {1. Reducción de Esfuerzos en el Centroide $G$};
                % Cilindro
                \draw[thick, fill=blue!6] (0, 0) -- (3.6, 0) arc (-90:90:0.45 and 0.9) -- (0, 1.8) arc (90:270:0.45 and 0.9);
                \draw[thick, fill=blue!15] (3.6, 0.9) ellipse (0.45 and 0.9);
                % Apoyo articulado izquierdo
                \draw[thick] (0.3, 0) -- (0.05, -0.45) -- (0.55, -0.45) -- cycle;
                \draw[thick] (-0.1, -0.45) -- (0.7, -0.45);
                % Fuerza exterior F
                \draw[->, line width=1.6pt, red!85!black] (0.9, 2.4) -- (1.3, 1.8) node[pos=0.0, left=2pt, font=\footnotesize\bfseries] {$\vec{F}$};
                % Ejes locales en G(3.6, 0.9)
                \coordinate (G) at (3.6, 0.9);
                \draw[->, thick, blue!85!black] (G) -- ++(2.2, 0) node[right, font=\footnotesize\bfseries] {$X$};
                \draw[->, thick, blue!85!black] (G) -- ++(0, 1.45) node[right=2pt, font=\footnotesize\bfseries] {$Y$};
                \draw[->, thick, blue!85!black] (G) -- ++(-0.85, -0.85) node[below left=-1pt, font=\footnotesize\bfseries] {$Z$};
                % Vectores R y M
                \draw[->, line width=1.8pt, black] (G) -- ++(1.7, 0.55) node[right, font=\footnotesize\bfseries] {$\vec{R}$};
                \draw[->, line width=1.5pt, purple!85!black] (G) -- ++(1.0, 0.85) node[above, font=\footnotesize\bfseries] {$\vec{M}$};
                \filldraw[black] (G) circle (2pt) node[below right=1pt, font=\scriptsize\bfseries] {$G$};
            \end{scope}

            % 2. Comparativa Tracción/Compresión vs Flexión
            \begin{scope}[shift={(7.5, 2.2)}]
                \node[above, font=\small\bfseries, blue!85!black] at (2.6, 2.75) {2. Solicitación de Flexión Plana en Vigas};
                % Viga prismática 3D bajo cargas transversales
                \draw[thick, fill=orange!10] (0, 0.3) rectangle (4.6, 1.3);
                \draw[thick, fill=orange!18] (0, 1.3) -- (0.5, 1.6) -- (5.1, 1.6) -- (4.6, 1.3) -- cycle;
                \draw[thick, fill=orange!22] (4.6, 0.3) -- (5.1, 0.6) -- (5.1, 1.6) -- (4.6, 1.3) -- cycle;
                \draw[dashdotted, gray!80, line width=0.8pt] (-0.3, 0.8) -- (5.4, 0.8);
                % Apoyo móvil izquierdo y fijo derecho
                \draw[thick] (0.4, 0.3) -- (0.2, -0.05) -- (0.6, -0.05) -- cycle;
                \draw[thick] (0.1, -0.12) -- (0.7, -0.12);
                \draw[thick] (5.1, 0.9) -- (5.45, 0.7) -- (5.45, 1.1) -- cycle;
                % Cargas transversales verticales
                \foreach \xx in {1.6, 2.5, 3.4} {
                        \draw[->, line width=1.5pt, red!85!black] (\xx, 2.15) -- (\xx, 1.45);
                    }
                \node[above=1pt, font=\scriptsize\bfseries, red!85!black] at (2.5, 2.15) {Cargas transversales $\perp X$};
            \end{scope}
        \end{tikzpicture}%
    }
    \caption{Reducción de esfuerzos internos en el centro de gravedad $G$ de la sección recta y esquema de una viga solicitada a flexión plana.}
    \label{fig:introduccion_flexion_esfuerzos}
\end{figure}

Contrastando el estado tensional con el capítulo anterior:
\begin{itemize}
    \item \textbf{Tracción / Compresión Uniaxial (Tema 4):} En la sección actúa únicamente el esfuerzo axil $N \neq 0$, generando tensiones normales uniformes $\sigma_n = \frac{N}{\Omega}$ y alargamientos/acortamientos longitudinales $\Delta l = \int_A^B \frac{N}{E\Omega}\,dx$.
    \item \textbf{Flexión Plana Directa (Tema 5):} Las cargas exteriores actúan contenidas en el plano de simetría vertical $XY$ y perpendicularmente al eje de la viga. En consecuencia, en la sección transversal aparecen como esfuerzos determinantes el \textbf{momento flector $M_z$} y, en general, el \textbf{esfuerzo cortante $Q_y$}:
          \begin{enumerate}
              \item \textbf{Flexión Pura:} En el tramo considerado actúa exclusivamente momento flector constante ($M_z \neq 0$, $Q_y = 0$, $N = 0$).
              \item \textbf{Flexión Simple:} En la sección coexisten simultáneamente el momento flector y el esfuerzo cortante ($M_z \neq 0$, $Q_y \neq 0$, $N = 0$).
          \end{enumerate}
\end{itemize}

%=============================================================================
\section{Análisis de Tensiones en Flexión Pura: Tensiones Normales y Fórmula de Navier}

\subsection{Modelo Clásico de Flexión Pura e Hipótesis de Partida}

El prototipo clásico de \textbf{flexión pura} se presenta en el tramo central $BC$ de una viga biapoyada simétrica de luz $L$ sometida a dos cargas puntuales idénticas $F$ situadas a igual distancia $a$ de los apoyos extremos $A$ y $D$:
\begin{itemize}
    \item En los tramos laterales $AB$ y $CD$ existe esfuerzo cortante ($Q_y = +F$ y $Q_y = -F$) y momento flector linealmente variable (flexión simple).
    \item En el \textbf{tramo central $BC$}, el esfuerzo axil y el esfuerzo cortante son idénticamente nulos ($N = 0$, $Q_y = 0$), mientras que el momento flector se mantiene estrictamente constante e igual a $\mathbf{M_z = F \cdot a}$ (\textbf{flexión pura}).
\end{itemize}

\begin{figure}[H]
    \centering
    \adjustbox{max width=0.92\textwidth}{%
        \begin{tikzpicture}[>=Latex, scale=0.92]
            % Columna izquierda: Viga biapoyada y diagramas N, Qy, Mz
            \begin{scope}[shift={(0,0)}]
                \node[above, font=\small\bfseries, red!85!black] at (3.5, 6.6) {ESQUEMA DE FLEXIÓN PURA (TRAMO CENTRAL $BC$)};

                % Banda sombreada del tramo BC (x de 2.0 a 5.0)
                \fill[red!8] (2.0, 0.0) rectangle (5.0, 5.8);
                \draw[dashed, gray!70] (2.0, 0.0) -- (2.0, 5.8);
                \draw[dashed, gray!70] (5.0, 0.0) -- (5.0, 5.8);

                % Viga en y = 5.4
                \draw[->, thick] (0, 5.4) -- (0, 6.4) node[right, font=\footnotesize] {$y$};
                \draw[->, thick] (0, 5.4) -- (7.6, 5.4) node[above, font=\footnotesize] {$x$};
                \draw[line width=2.4pt, black] (0, 5.4) -- (7.0, 5.4);
                % Apoyos A y D
                \draw[thick] (0, 5.4) -- (-0.25, 4.95) -- (0.25, 4.95) -- cycle;
                \draw[thick] (-0.4, 4.95) -- (0.4, 4.95);
                \draw[thick] (7.0, 5.4) -- (6.75, 4.95) -- (7.25, 4.95) -- cycle;
                \draw[thick] (6.6, 4.88) -- (7.4, 4.88);
                % Nodos y cargas F
                \node[above left=1pt, font=\footnotesize\bfseries] at (0, 5.4) {$A$};
                \node[above left=1pt, font=\footnotesize\bfseries] at (2.0, 5.4) {$B$};
                \node[above right=1pt, font=\footnotesize\bfseries] at (5.0, 5.4) {$C$};
                \node[above right=1pt, font=\footnotesize\bfseries] at (7.0, 5.4) {$D$};
                \draw[->, line width=1.6pt, red!85!black] (2.0, 6.3) node[above=1pt, font=\footnotesize\bfseries] {$F$} -- (2.0, 5.45);
                \draw[->, line width=1.6pt, red!85!black] (5.0, 6.3) node[above=1pt, font=\footnotesize\bfseries] {$F$} -- (5.0, 5.45);
                % Cotas a
                \draw[<->, line width=0.7pt] (0, 4.6) -- node[below=1pt, font=\footnotesize] {$a$} (2.0, 4.6);
                \draw[<->, line width=0.7pt] (5.0, 4.6) -- node[below=1pt, font=\footnotesize] {$a$} (7.0, 4.6);

                % Diagrama N (y = 3.6)
                \draw[thick] (0, 3.6) -- (7.0, 3.6);
                \node[left=3pt, font=\footnotesize\bfseries] at (0, 3.6) {$N$};
                \node[above=1pt, font=\footnotesize\bfseries, blue!85!black] at (3.5, 3.6) {$N = 0$};

                % Diagrama Qy (y = 2.2)
                \draw[thick] (0, 2.2) -- (7.0, 2.2);
                \node[left=3pt, font=\footnotesize\bfseries] at (0, 2.2) {$Q_y$};
                \draw[thick, fill=blue!12] (0, 2.2) rectangle (2.0, 2.85);
                \node[above=1pt, font=\footnotesize\bfseries] at (1.0, 2.85) {$+F$};
                \node[above=1pt, font=\footnotesize\bfseries, blue!85!black] at (3.5, 2.2) {$Q_y = 0$};
                \draw[thick, fill=red!12] (5.0, 2.2) rectangle (7.0, 1.55);
                \node[below=1pt, font=\footnotesize\bfseries] at (6.0, 1.55) {$-F$};

                % Diagrama Mz (y = 0.9)
                \draw[thick] (0, 0.9) -- (7.0, 0.9);
                \node[left=3pt, font=\footnotesize\bfseries] at (0, 0.9) {$M_z$};
                \draw[thick, purple!85!black, fill=purple!12] (0, 0.9) -- (2.0, 0.0) -- (5.0, 0.0) -- (7.0, 0.9) -- cycle;
                \node[below=2pt, font=\footnotesize\bfseries, purple!85!black] at (3.5, 0.0) {$M_z = F \cdot a = \text{cte}$};
            \end{scope}

            % Columna derecha: Deformación de rebanada según Hipótesis de Navier-Bernoulli
            \begin{scope}[shift={(9.0, 1.2)}]
                \node[above, font=\small\bfseries, blue!85!black] at (2.5, 4.6) {Hipótesis de Secciones Planas (Bernoulli)};
                % Rebanada sin deformar
                \draw[thick, fill=gray!10] (0.5, 2.8) rectangle (4.5, 4.0);
                \foreach \xx in {1.3, 2.1, 2.9, 3.7} {
                        \draw[thin, gray!70] (\xx, 2.8) -- (\xx, 4.0);
                    }
                \draw[dashed, blue!80!black] (0.3, 3.4) -- (4.7, 3.4);
                \node[below=1pt, font=\scriptsize\bfseries] at (2.5, 2.8) {Configuración indeformada ($\gamma_{xy} = 0$)};

                % Rebanada deformada por flexión pura
                \draw[thick, blue!85!black] (0.5, 1.6) to[bend right=18] (4.5, 1.6);
                \draw[thick, blue!85!black] (0.5, 0.5) to[bend right=18] (4.5, 0.5);
                \draw[dashed, red!85!black, line width=1.0pt] (0.5, 1.05) to[bend right=18] (4.5, 1.05);
                \draw[thick, blue!85!black] (0.5, 1.6) -- (0.8, 0.52);
                \draw[thick, blue!85!black] (4.5, 1.6) -- (4.2, 0.52);
                \draw[->, thick, purple!85!black] (0.2, 0.6) arc (210:140:0.8) node[left=1pt, font=\footnotesize\bfseries] {$M_z$};
                \draw[->, thick, purple!85!black] (4.8, 0.6) arc (-30:40:0.8) node[right=1pt, font=\footnotesize\bfseries] {$M_z$};
                \node[below=2pt, font=\scriptsize\bfseries, align=center] at (2.5, 0.2) {Secciones planas y $\perp$ directriz deformada\\$\implies \gamma = 0 \implies \tau = 0$ (solo $\sigma_{nx}$)};
            \end{scope}
        \end{tikzpicture}%
    }
    \caption{Caracterización de la flexión pura ($N=0, Q_y=0, M_z=\text{cte}$) y conservación de la planaridad de las secciones rectas según la hipótesis de Navier-Bernoulli.}
    \label{fig:flexion_pura_hipotesis}
\end{figure}

Las \textbf{hipótesis fundamentales de partida} para deducir la ley tensional en flexión pura son:
\begin{enumerate}
    \item \textbf{Comportamiento elástico lineal (Ley de Hooke):} El material obedece la relación constitutiva lineal con estado tensional uniaxial dominante ($\sigma_{ny} = \sigma_{nz} = 0$):
          \begin{equation}
              \varepsilon_{nx} = \frac{1}{E}\left(\sigma_{nx} - \nu(\cancel{\sigma_{ny} + \sigma_{nz}})\right) = \frac{\sigma_{nx}}{E} \implies \mathbf{\sigma_{nx} = E \cdot \varepsilon_{nx}}, \qquad \tau_{xy} = G \cdot \gamma_{xy}
          \end{equation}
    \item \textbf{Hipótesis cinemática de Navier-Bernoulli:} \textit{Las secciones transversales originariamente planas y perpendiculares a la directriz de la viga se mantienen planas y perpendiculares a la directriz deformada tras aplicar las cargas} (válida para vigas esbeltas con relación luz/canto $L \ge 10 h$).
          Al no producirse distorsión angular ($\gamma_{xy} = 0$), las tensiones tangenciales son nulas ($\tau_{xy} = 0$) y \textbf{únicamente existen tensiones normales $\sigma_{nx}$}.
    \item \textbf{Principio de rigidez relativa y pequeñas deformaciones:} Los desplazamientos $\delta$ son suficientemente pequeños frente a las dimensiones de la estructura como para plantear el equilibrio en la geometría sin deformar.
    \item \textbf{Principio de superposición de efectos:} Consecuencia directa de la linealidad mecánica (Ley de Hooke) y geométrica (rigidez relativa).
    \item \textbf{Simetría geométrica y de carga:} La sección transversal es simétrica respecto al eje vertical $Y$, y las cargas están contenidas en el plano de simetría $XY$ (flexión recta o simétrica).
\end{enumerate}

%=============================================================================
\subsection{Deducción de la Fórmula de Navier}

\subsubsection{1. Relación Cinemática Tensión--Curvatura}

Consideremos una rebanada elemental de viga de longitud inicial $dx$ sometida a un momento flector positivo $M_z > 0$ (que comprime las fibras superiores $y > 0$ y tracciona las fibras inferiores $y < 0$).

Existe una superficie longitudinal cuyas fibras no sufren ni alargamiento ni acortamiento ($\varepsilon_x = 0$), denominada \textbf{superficie neutra}. Su intersección con el plano de la sección transversal define el \textbf{eje neutro (EN)}, y su intersección con el plano de simetría longitudinal $XY$ se denomina \textbf{fibra neutra (FN)}, de radio de curvatura $\rho$.

\begin{figure}[H]
    \centering
    \adjustbox{max width=0.92\textwidth}{%
        \begin{tikzpicture}[>=Latex, scale=1.0]
            % 1. Rebanada deformada izquierda
            \begin{scope}[shift={(0,0)}]
                \node[above, font=\small\bfseries, blue!85!black] at (2.8, 3.8) {1. Deformación de las Fibras Longitudinales};
                % Fibra superior, neutra e inferior
                \draw[thick, blue!85!black] (0.4, 2.6) to[bend right=22] (5.2, 2.6);
                \draw[dashed, red!85!black, line width=1.1pt] (0.2, 1.8) to[bend right=22] (5.4, 1.8);
                \draw[thick, blue!85!black] (0.0, 1.0) to[bend right=22] (5.6, 1.0);
                \draw[thick, blue!85!black] (0.4, 2.6) -- (0.0, 1.0);
                \draw[thick, blue!85!black] (5.2, 2.6) -- (5.6, 1.0);
                % Rebanada central AB
                \draw[thin, gray!80] (2.4, 0.65) -- (2.5, 2.8);
                \draw[thin, gray!80] (3.2, 0.65) -- (3.1, 2.8);
                \node[below left=-1pt, font=\scriptsize\bfseries] at (2.45, 1.45) {$A$};
                \node[below right=-1pt, font=\scriptsize\bfseries] at (3.15, 1.45) {$B$};
                % Momentos Mz
                \draw[->, line width=1.4pt, purple!85!black] (-0.45, 0.9) arc (225:145:0.85) node[above left=1pt, font=\footnotesize\bfseries] {$M_z$};
                \draw[->, line width=1.4pt, purple!85!black] (6.05, 0.9) arc (-45:35:0.85) node[above right=1pt, font=\footnotesize\bfseries] {$M_z$};
                % Etiquetas
                \node[above=2pt, font=\scriptsize\bfseries, blue!85!black] at (2.8, 2.65) {Fibras superiores ($y > 0$, comprimidas)};
                \node[left=4pt, font=\scriptsize\bfseries, red!85!black] at (-0.65, 1.8) {Fibra neutra ($\varepsilon_x = 0$)};
                \node[below=2pt, font=\scriptsize\bfseries, green!55!black] at (2.8, 0.38) {Fibras inferiores ($y < 0$, traccionadas)};
            \end{scope}

            % 2. Semejanza geométrica de triángulos con centro de curvatura O
            \begin{scope}[shift={(8.2, 0)}]
                \node[above, font=\small\bfseries, blue!85!black] at (2.2, 3.8) {2. Geometría del Radio de Curvatura $\rho$};
                \coordinate (O) at (2.2, 3.5);
                \filldraw[black] (O) circle (2pt) node[above=2pt, font=\footnotesize\bfseries] {$O$};
                % Líneas radiales desde O
                \draw[thick, gray!80] (O) -- (1.0, 0.3);
                \draw[thick, gray!80] (O) -- (3.4, 0.3);
                \draw[dashdotted, gray!70] (O) -- (2.2, 0.2);
                % Fibra neutra AB a radio rho
                \draw[dashed, red!85!black, line width=1.1pt] (1.26, 1.0) to[bend right=14] (3.14, 1.0);
                \node[below left=-1pt, font=\scriptsize\bfseries] at (1.26, 1.0) {$A$};
                \node[below right=-1pt, font=\scriptsize\bfseries] at (3.14, 1.0) {$B$};
                % Fibra a altura y > 0
                \draw[thick, blue!85!black] (1.56, 1.8) to[bend right=12] (2.84, 1.8);
                % Paralela desde B
                \draw[dashed, purple!85!black, line width=0.9pt] (3.14, 1.0) -- (2.54, 2.6);
                \node[above left=-1pt, font=\scriptsize\bfseries] at (2.68, 1.85) {$M$};
                \node[above right=-1pt, font=\scriptsize\bfseries] at (2.90, 1.85) {$N$};
                % Cotas rho, y, dx
                \draw[<->, line width=0.7pt] (1.75, 1.0) -- node[left=1pt, font=\footnotesize\bfseries] {$\rho$} (1.75, 3.5);
                \draw[<->, line width=0.7pt] (3.6, 1.0) -- node[right=1pt, font=\footnotesize\bfseries] {$y > 0$} (3.6, 1.8);
                \draw[<->, line width=0.7pt] (1.26, 0.45) -- node[below=1pt, font=\footnotesize\bfseries] {$dx$} (3.14, 0.45);
            \end{scope}
        \end{tikzpicture}%
    }
    \caption{Relación geométrica entre la deformación longitudinal $\varepsilon_x$, la ordenada $y$ respecto a la fibra neutra y el radio de curvatura $\rho$.}
    \label{fig:relacion_tension_curvatura}
\end{figure}

Por semejanza geométrica entre el sector circular $OAB$ (de radio $\rho$ y arco $AB = dx$) y el triángulo diferencial $BMN$ (de cateto radial $BN = y$ y acortamiento de fibra $MN = -\Delta dx$):
\begin{equation}
    \frac{AB}{OA} = \frac{MN}{BN} \implies \frac{dx / 2}{\rho} = \frac{-\Delta dx / 2}{y} \implies \varepsilon_x = \frac{\Delta dx}{dx} = -\frac{y}{\rho}
\end{equation}
Sustituyendo en la Ley de Hooke ($\sigma_{nx} = E \cdot \varepsilon_x$), obtenemos la \textbf{relación tensión--curvatura}:
\begin{equation}
    \boxed{\sigma_{nx} = -\frac{E \cdot y}{\rho}}
    \label{eq:relacion_tension_curvatura}
\end{equation}

\subsubsection{2. Posición de la Fibra Neutra ($N = 0$)}

Imponiendo que en flexión pura el esfuerzo axil resultante sobre la sección es nulo ($N = 0$):
\begin{equation}
    N = \iint_{\Omega} \sigma_{nx} \, d\Omega = \iint_{\Omega} -\frac{E \cdot y}{\rho} \, d\Omega = -\frac{E}{\rho} \iint_{\Omega} y \, d\Omega = -\frac{E}{\rho} \cdot S_z = 0
\end{equation}
Como la rigidez elástica $E$ y la curvatura $1/\rho$ no son nulas, se deduce que el momento estático respecto al eje neutro debe anularse:
\begin{equation}
    \mathbf{S_z = \iint_{\Omega} y \, d\Omega = 0 \implies \text{\textbf{El eje neutro ($z$) pasa obligatoriamente por el centro de gravedad $G$ de la sección.}}}
\end{equation}

\subsubsection{3. Relación Momento--Curvatura y Fórmula de Navier ($M_z \neq 0$)}

Sustituyendo la expresión tensional~\eqref{eq:relacion_tension_curvatura} en la ecuación de equivalencia del momento flector $M_z$:
\begin{equation}
    M_z = \iint_{\Omega} -\sigma_{nx} \cdot y \, d\Omega = \iint_{\Omega} -\left(-\frac{E \cdot y}{\rho}\right) y \, d\Omega = \frac{E}{\rho} \iint_{\Omega} y^2 \, d\Omega = \frac{E \cdot I_z}{\rho}
\end{equation}
De aquí se obtiene la ecuación fundamental de la \textbf{curvatura elástica $\kappa$}:
\begin{equation}
    \boxed{\frac{1}{\rho} = \kappa = \frac{M_z}{E \cdot I_z}}
    \label{eq:curvatura_flexion}
\end{equation}
Finalmente, despejando $\frac{1}{\rho} = \frac{M_z}{E I_z}$ e introduciéndolo en $\sigma_{nx} = -\frac{E y}{\rho}$, llegamos a la célebre \textbf{Fórmula de Navier}:
\begin{equation}
    \boxed{\sigma_{nx}(y) = -\frac{M_z}{I_z} \cdot y}
    \label{eq:formula_navier}
\end{equation}

\begin{tcolorbox}[colback=green!5!white, colframe=myg!85!black, title=\textbf{Interpretación Física de la Fórmula de Navier y Distribución Lineal}]
    En flexión pura, las tensiones normales $\sigma_{nx}(y)$ varían \textbf{linealmente con la distancia $y$ a la fibra neutra}:
    \begin{itemize}
        \item En la \textbf{fibra neutra ($y = 0$, que pasa por $G$):} la tensión normal es nula ($\sigma_{nx} = 0$).
        \item Para un \textbf{momento flector positivo ($M_z > 0$, carita sonriente $\smile$):}
              \begin{itemize}
                  \item Las \textbf{fibras superiores ($y > 0$)} están sometidas a \textbf{compresión ($\sigma_{nx} < 0$)}, alcanzando su valor extremo en la coronación superior ($y = +h_1$): $\sigma_{\text{sup}} = -\frac{M_z \cdot h_1}{I_z}$.
                  \item Las \textbf{fibras inferiores ($y < 0$)} están sometidas a \textbf{tracción ($\sigma_{nx} > 0$)}, alcanzando su valor extremo en el borde inferior ($y = -h_2$): $\sigma_{\text{inf}} = +\frac{M_z \cdot h_2}{I_z}$.
              \end{itemize}
        \item Para un \textbf{momento flector negativo ($M_z < 0$, carita triste $\frown$):} los signos se invierten (tracción arriba en $y > 0$ y compresión abajo en $y < 0$).
    \end{itemize}
\end{tcolorbox}

\begin{figure}[H]
    \centering
    \adjustbox{max width=0.88\textwidth}{%
        \begin{tikzpicture}[>=Latex, scale=1.0]
            % 1. Sección transversal izquierda
            \begin{scope}[shift={(0,0)}]
                \fill[blue!10] (0, 0) rectangle (2.2, 3.2);
                \draw[thick, blue!85!black] (0, 0) rectangle (2.2, 3.2);
                \draw[dashed, red!85!black, line width=1.1pt] (-0.5, 1.4) -- (5.5, 1.4);
                \node[below=2pt, font=\scriptsize\bfseries, red!85!black] at (1.1, 1.4) {Fibra neutra};
                \draw[->, thick, black] (0.8, 1.4) -- (0.8, 2.9) node[midway, right=2pt, font=\footnotesize\bfseries] {$y > 0$};
                % Cotas h1 y h2
                \draw[<->, line width=0.7pt] (-0.35, 1.4) -- node[left=1pt, font=\footnotesize] {$h_1$} (-0.35, 3.2);
                \draw[<->, line width=0.7pt] (-0.35, 0.0) -- node[left=1pt, font=\footnotesize] {$h_2$} (-0.35, 1.4);

                % Diagrama bitriangular de tensiones a la derecha de la sección (x = 3.2)
                \draw[thick] (3.2, 0.0) -- (3.2, 3.2);
                \fill[red!12] (3.2, 1.4) -- (3.2, 3.2) -- (4.6, 3.2) -- cycle;
                \fill[green!12] (3.2, 1.4) -- (3.2, 0.0) -- (4.3, 0.0) -- cycle;
                \draw[thick, blue!85!black] (4.3, 0.0) -- (3.2, 1.4) -- (4.6, 3.2);
                % Flechas de compresión (arriba, hacia la sección) y tracción (abajo, salientes)
                \draw[<-, thick, red!85!black] (3.22, 3.2) -- (4.6, 3.2) node[midway, above=2pt, font=\footnotesize\bfseries] {$\sigma_{\text{máx},C} = -\dfrac{M_z h_1}{I_z}$};
                \draw[<-, thin, red!85!black] (3.22, 2.6) -- (4.13, 2.6);
                \draw[<-, thin, red!85!black] (3.22, 2.0) -- (3.67, 2.0);
                \draw[->, thin, green!60!black] (3.2, 0.93) -- (3.57, 0.93);
                \draw[->, thin, green!60!black] (3.2, 0.47) -- (3.93, 0.47);
                \draw[->, thick, green!60!black] (3.2, 0.0) -- (4.3, 0.0) node[midway, below=2pt, font=\footnotesize\bfseries] {$\sigma_{\text{máx},T} = +\dfrac{M_z h_2}{I_z}$};
                % Momento Mz
                \draw[->, line width=1.4pt, purple!85!black] (4.9, 1.0) arc (-40:40:0.9) node[midway, right=2pt, font=\footnotesize\bfseries] {$M_z > 0$};
            \end{scope}

            % 2. Esquema 3D del prisma de tensiones
            \begin{scope}[shift={(8.2, 0.2)}]
                \draw[thick, blue!85!black] (0, 0) -- (0, 3.0) -- (1.4, 3.6) -- (1.4, 0.6) -- cycle;
                \draw[dashed, red!85!black, line width=1.1pt] (0, 1.5) -- (1.4, 2.1);
                \node[left=3pt, font=\scriptsize\bfseries, red!85!black] at (0, 1.5) {Eje neutro};
                % Cuñas 3D
                \draw[thick, red!80!black, fill=red!10, fill opacity=0.7] (0, 1.5) -- (0, 3.0) -- (-1.1, 3.0) -- cycle;
                \draw[thick, red!80!black, fill=red!10, fill opacity=0.7] (1.4, 2.1) -- (1.4, 3.6) -- (0.3, 3.6) -- cycle;
                \draw[thick, red!80!black] (-1.1, 3.0) -- (0.3, 3.6);
                \draw[thick, green!60!black, fill=green!10, fill opacity=0.7] (0, 1.5) -- (0, 0) -- (1.1, 0) -- cycle;
                \draw[thick, green!60!black, fill=green!10, fill opacity=0.7] (1.4, 2.1) -- (1.4, 0.6) -- (2.5, 0.6) -- cycle;
                \draw[thick, green!60!black] (1.1, 0) -- (2.5, 0.6);
                \draw[->, line width=1.4pt, purple!85!black] (2.8, 1.2) arc (-40:40:0.9) node[midway, right=2pt, font=\footnotesize\bfseries] {$M_z$};
            \end{scope}
        \end{tikzpicture}%
    }
    \caption{Distribución lineal bitriangular de tensiones normales $\sigma_{nx}(y)$ en flexión pura según la ley de Navier.}
    \label{fig:distribucion_navier_3d}
\end{figure}

%=============================================================================
\section{Módulo Resistente y Formas Más Adecuadas de la Sección Transversal}

En el diseño y comprobación de vigas a flexión interesa conocer las \textbf{tensiones normales máximas} que se alcanzan en las fibras más alejadas de la fibra neutra ($y_{\text{máx,sup}} = h_1$ e $|y_{\text{máx,inf}}| = h_2$). Operando en valor absoluto:
\begin{equation}
    \sigma_{nx,\text{máx}} = \frac{|M_z| \cdot y_{\text{máx}}}{I_z} = \frac{|M_z|}{\left(\dfrac{I_z}{y_{\text{máx}}}\right)} = \frac{|M_z|}{W_z}
\end{equation}
Se define el \textbf{Módulo Resistente elástico de la sección ($W_z$)} como el cociente entre el momento de inercia respecto al eje neutro y la distancia a la fibra extrema:
\begin{equation}
    \boxed{W_{z,\text{sup}} = \frac{I_z}{y_{\text{máx,sup}}}, \qquad W_{z,\text{inf}} = \frac{I_z}{|y_{\text{máx,inf}}|}}
    \label{eq:modulo_resistente}
\end{equation}
Sus unidades son de longitud al cubo ($\text{mm}^3$, $\text{cm}^3$ o $\text{m}^3$). En secciones simétricas respecto al eje $z$ ($h_1 = h_2 = h/2$), ambos módulos coinciden ($W_z = \frac{I_z}{h/2}$) y su valor se encuentra directamente tabulado en los prontuarios de perfiles metálicos laminados y conformados.

\subsection{Eficiencia Estructural y Perfiles Más Usados en Edificación}

Dado que la tensión normal crece linealmente con la distancia $y$ al eje neutro ($\sigma \propto y$) y el momento de inercia crece con el cuadrado de dicha distancia ($I_z = \iint y^2 d\Omega$), el material situado cerca de la fibra neutra apenas aporta resistencia a flexión.

Por consiguiente, las \textbf{formas más adecuadas para resistir flexión} son aquellas que \textbf{concentran la mayor parte del área lejos del eje neutro (en las alas superior e inferior)}, conectándolas mediante un alma delgada:
\begin{itemize}
    \item \textbf{Vigas y viguetas (flexión dominante en un plano):} Perfiles en doble T esbeltos de tipo \textbf{IPE} e \textbf{IPN}, así como perfiles en U (\textbf{UPN}), que maximizan $I_z$ y $W_z$ con el mínimo peso propio.
    \item \textbf{Pilares y soportes (compresión combinada y flexión esviada):} Perfiles de alas anchas tipo \textbf{HEB, HEA y HEM}, que proporcionan una mayor inercia transversal $I_y$ frente al pandeo y la flexión en ambos planos.
\end{itemize}

%=============================================================================
\section{Ejercicios Resueltos de Comprobación Tensional y Eficiencia en Flexión}

\ex{Comprobación de Viga Armada en T bajo Momento Positivo y Negativo}{
    Dada la viga armada de acero estructural \textbf{S275} ($f_{yk} = 275\text{ N/mm}^2$) cuya sección transversal en T tiene las dimensiones indicadas en la figura:
    \begin{itemize}
        \item \textbf{Ala superior horizontal:} ancho $b_f = 250\text{ mm}$ y espesor $t_f = 8\text{ mm}$.
        \item \textbf{Alma vertical inferior:} altura $h_w = 300\text{ mm}$ y espesor $t_w = 6\text{ mm}$.
    \end{itemize}
    \textbf{Se pide:}
    \begin{enumerate}[label=\alph*)]
        \item Comprobar si la sección es apta para resistir un momento flector positivo de cálculo $M_d = +45\text{ kN}\cdot\text{m}$.
        \item Comprobar si la sección es apta para resistir un momento flector negativo de cálculo $M_d = -38\text{ kN}\cdot\text{m}$.
    \end{enumerate}

    \begin{center}
        \adjustbox{max width=0.68\textwidth}{%
            \begin{tikzpicture}[>=Latex, scale=0.92]
                % Eje z en la base
                \draw[->, line width=1.0pt] (0, 0) -- (5.8, 0) node[right, font=\bfseries] {$z$};
                % Alma vertical (6 mm x 300 mm -> representamos ancho 0.35 y alto 3.6)
                \fill[green!12] (2.325, 0) rectangle (2.675, 3.6);
                \draw[thick, green!60!black] (2.325, 0) rectangle (2.675, 3.6);
                % Ala superior (250 mm x 8 mm -> ancho 4.2, alto 0.35)
                \fill[blue!12] (0.4, 3.6) rectangle (4.6, 3.95);
                \draw[thick, blue!85!black] (0.4, 3.6) rectangle (4.6, 3.95);

                % Cotas
                \draw[<->, line width=0.7pt] (0.4, 4.3) -- node[above=1pt, font=\footnotesize\bfseries] {$250\text{ mm}$} (4.6, 4.3);
                \draw[<->, line width=0.7pt] (2.25, -0.45) -- node[below=1pt, font=\footnotesize\bfseries] {$6\text{ mm}$} (2.75, -0.45);
                \draw[<->, line width=0.7pt] (5.1, 0) -- node[right=2pt, font=\footnotesize\bfseries] {$300\text{ mm}$} (5.1, 3.6);
                \draw[<->, line width=0.7pt] (5.1, 3.6) -- node[right=2pt, font=\footnotesize\bfseries] {$8\text{ mm}$} (5.1, 3.95);
            \end{tikzpicture}%
        }
    \end{center}

    \tcblower

    \paso{1}{Propiedades Geométricas: Centro de Gravedad e Inercia Baricéntrica}

    \begin{itemize}
        \item \textbf{Coordenada horizontal $Z_G$:} Por simetría respecto al eje vertical del alma:
              \begin{equation}
                  \mathbf{Z_G = \frac{250}{2} = 125\text{ mm}}
              \end{equation}
        \item \textbf{Coordenada vertical $Y_G$ (medida desde el borde inferior del alma $y = 0$):}
              Con $\Omega_{\text{alma}} = 300 \times 6 = 1800\text{ mm}^2$ ($Y_{G,1} = 150\text{ mm}$) y $\Omega_{\text{ala}} = 250 \times 8 = 2000\text{ mm}^2$ ($Y_{G,2} = 300 + 4 = 304\text{ mm}$):
              \begin{equation}
                  Y_G = \frac{1800 \cdot 150 + 2000 \cdot 304}{1800 + 2000} = \frac{270\,000 + 608\,000}{3800} = \frac{878\,000}{3800} = \mathbf{231.05\text{ mm}}
              \end{equation}
              Por tanto, las distancias desde la fibra neutra ($Y_G = 231.05\text{ mm}$) hasta las fibras extremas inferior y superior son:
              \begin{equation}
                  y_{\text{inf}} = \mathbf{231.05\text{ mm}}, \qquad y_{\text{sup}} = 308 - 231.05 = \mathbf{76.95\text{ mm}}
              \end{equation}
        \item \textbf{Momento de Inercia $I_{z_G}$ (Teorema de Steiner):}
              \begin{align}
                  I_{z_G} & = \left[ \frac{6 \cdot 300^3}{12} + 1800 \cdot (231.05 - 150)^2 \right] + \left[ \frac{250 \cdot 8^3}{12} + 2000 \cdot (304 - 231.05)^2 \right] \notag \\
                          & = 25\,324\,384.5 + 10\,654\,081.7 = \mathbf{35.978467 \times 10^6\text{ mm}^4}
              \end{align}
    \end{itemize}

    \paso{2}{Apartado a) Comprobación para $M_d = +45\text{ kN}\cdot\text{m} = 45 \times 10^6\text{ N}\cdot\text{mm}$}

    Un momento flector positivo ($M_d > 0$) comprime las fibras situadas por encima de la fibra neutra ($y_{\text{sup}} = 76.95\text{ mm}$) y tracciona las fibras inferiores ($y_{\text{inf}} = 231.05\text{ mm}$):
    \begin{align}
        \sigma_{\text{máx},C} & = -\frac{M_d \cdot y_{\text{sup}}}{I_{z_G}} = -\frac{45 \times 10^6 \cdot 76.95}{35.978467 \times 10^6} \approx \mathbf{-96.25\text{ N/mm}^2} \approx \mathbf{-96\text{ N/mm}^2}    \\
        \sigma_{\text{máx},T} & = +\frac{M_d \cdot y_{\text{inf}}}{I_{z_G}} = +\frac{45 \times 10^6 \cdot 231.05}{35.978467 \times 10^6} \approx \mathbf{+288.98\text{ N/mm}^2} \approx \mathbf{+289\text{ N/mm}^2}
    \end{align}
    \textbf{Verificación resistente:} Como la tensión máxima de tracción en la fibra inferior alcanza $|\sigma_{\text{máx},T}| = 289\text{ N/mm}^2 > f_{yk} = 275\text{ N/mm}^2$, la sección \textbf{NO CUMPLE (NO es apta)} para resistir $M_d = +45\text{ kN}\cdot\text{m}$.

    \paso{3}{Apartado b) Comprobación para $M_d = -38\text{ kN}\cdot\text{m} = -38 \times 10^6\text{ N}\cdot\text{mm}$}

    Un momento flector negativo ($M_d < 0$) tracciona las fibras superiores ($y_{\text{sup}} = 76.95\text{ mm}$) y comprime las fibras inferiores del alma ($y_{\text{inf}} = 231.05\text{ mm}$):
    \begin{align}
        \sigma_{\text{máx},T} & = +\frac{|M_d| \cdot y_{\text{sup}}}{I_{z_G}} = +\frac{38 \times 10^6 \cdot 76.95}{35.978467 \times 10^6} \approx \mathbf{+81.27\text{ N/mm}^2} \approx \mathbf{+81\text{ N/mm}^2}    \\
        \sigma_{\text{máx},C} & = -\frac{|M_d| \cdot y_{\text{inf}}}{I_{z_G}} = -\frac{38 \times 10^6 \cdot 231.05}{35.978467 \times 10^6} \approx \mathbf{-244.03\text{ N/mm}^2} \approx \mathbf{-244\text{ N/mm}^2}
    \end{align}
    \textbf{Verificación resistente:} La tensión máxima absoluta es $|\sigma_{\text{máx},C}| = 244\text{ N/mm}^2 \le f_{yk} = 275\text{ N/mm}^2$ (y también menor que $f_{yd} = 275/1.05 = 261.9\text{ N/mm}^2$). Por tanto, la sección \textbf{SÍ CUMPLE (es apta)} para resistir $M_d = -38\text{ kN}\cdot\text{m}$.

    \begin{center}
        \adjustbox{max width=0.92\textwidth}{%
            \begin{tikzpicture}[>=Latex, scale=0.88]
                % Diagrama a) Md = +45 kNm
                \begin{scope}[shift={(0,0)}]
                    \node[above, font=\small\bfseries, blue!85!black] at (2.2, 4.6) {a) $M_d = +45\text{ kN}\cdot\text{m}$ (NO apta)};
                    % Esquema perfil T a la izquierda
                    \draw[thick] (0, 3.6) -- (1.8, 3.6);
                    \draw[thick] (0.9, 0) -- (0.9, 3.6);
                    % Eje neutro en y = 2.7
                    \draw[dashdotted, gray!80, line width=0.9pt] (-0.2, 2.7) -- (4.5, 2.7) node[right, font=\footnotesize\bfseries] {$x_G$};
                    \draw[<->, line width=0.7pt] (0.2, 0) -- node[left=1pt, font=\scriptsize] {$231\text{ mm}$} (0.2, 2.7);
                    % Ley de tensiones en x = 2.6
                    \draw[thick] (2.6, 0) -- (2.6, 3.6);
                    \fill[red!12] (2.6, 2.7) -- (2.6, 3.6) -- (2.0, 3.6) -- cycle;
                    \fill[green!12] (2.6, 2.7) -- (2.6, 0) -- (1.1, 0) -- cycle;
                    \draw[thick, blue!85!black] (2.0, 3.6) -- (2.6, 2.7) -- (1.1, 0);
                    \draw[<-, thick, red!85!black] (2.0, 3.6) -- (2.6, 3.6);
                    \draw[->, thick, green!60!black] (1.1, 0) -- (2.6, 0);
                    \node[above=2pt, font=\scriptsize\bfseries, red!85!black] at (2.3, 3.6) {$\sigma_{\text{máx},C} = -96\text{ N/mm}^2$};
                    \node[below=2pt, font=\scriptsize\bfseries, red!85!black] at (2.0, 0) {$\sigma_{\text{máx},T} = +289\text{ N/mm}^2 > 275$};
                    \draw[->, line width=1.3pt, purple!85!black] (3.3, 1.4) arc (-45:45:1.1) node[pos=0.2, right=2pt, font=\scriptsize\bfseries] {$45\text{ kN}\cdot\text{m}$};
                \end{scope}

                % Diagrama b) Md = -38 kNm
                \begin{scope}[shift={(7.5,0)}]
                    \node[above, font=\small\bfseries, green!55!black] at (2.2, 4.6) {b) $M_d = -38\text{ kN}\cdot\text{m}$ (SÍ apta)};
                    \draw[thick] (0, 3.6) -- (1.8, 3.6);
                    \draw[thick] (0.9, 0) -- (0.9, 3.6);
                    \draw[dashdotted, gray!80, line width=0.9pt] (-0.2, 2.7) -- (4.5, 2.7) node[right, font=\footnotesize\bfseries] {$x_G$};
                    % Ley de tensiones en x = 2.6
                    \draw[thick] (2.6, 0) -- (2.6, 3.6);
                    \fill[green!12] (2.6, 2.7) -- (2.6, 3.6) -- (2.1, 3.6) -- cycle;
                    \fill[red!12] (2.6, 2.7) -- (2.6, 0) -- (1.3, 0) -- cycle;
                    \draw[thick, blue!85!black] (2.1, 3.6) -- (2.6, 2.7) -- (1.3, 0);
                    \draw[->, thick, green!60!black] (2.1, 3.6) -- (2.6, 3.6);
                    \draw[<-, thick, red!85!black] (1.3, 0) -- (2.6, 0);
                    \node[above=2pt, font=\scriptsize\bfseries, green!55!black] at (2.4, 3.6) {$\sigma_{\text{máx},T} = +81\text{ N/mm}^2$};
                    \node[below=2pt, font=\scriptsize\bfseries, blue!85!black] at (2.0, 0) {$\sigma_{\text{máx},C} = -244\text{ N/mm}^2 \le 275$};
                    \draw[<-, line width=1.3pt, purple!85!black] (3.3, 1.4) arc (-45:45:1.1) node[pos=0.2, right=2pt, font=\scriptsize\bfseries] {$38\text{ kN}\cdot\text{m}$};
                \end{scope}
            \end{tikzpicture}%
        }
    \end{center}
}

\ex{Comparativa de Capacidad Resistente a Tracción y Flexión de Tres Perfiles de Igual Área}{
    Se dispone de 3 perfiles metálicos de acero \textbf{S275} ($f_{yk} = 275\text{ N/mm}^2$) con idéntica altura total $H = 50\text{ cm}$ y distinta forma de sección transversal (cotas en $\text{cm}$, coeficientes parciales de seguridad $\gamma_s = 1.05$ y $\gamma_f = 1.50$):
    \begin{enumerate}
        \item \textbf{Sección 1 (Perfil en Doble T simétrico):} Dos alas horizontales de $50 \times 10\text{ cm}^2$ (sobresaliendo $20\text{ cm}$ a cada lado del alma) y un alma central de $10 \times 30\text{ cm}^2$.
        \item \textbf{Sección 2 (Perfil en Cruz simétrico):} Un nervio vertical central de $10 \times 50\text{ cm}^2$ y dos brazos horizontales centrados de $20 \times 20\text{ cm}^2$ cada uno (en la franja central de altura $20\text{ cm}$, dejando $15\text{ cm}$ libres arriba y abajo).
        \item \textbf{Sección 3 (Perfil Rectangular macizo):} Ancho $b = 26\text{ cm}$ y altura $h = 50\text{ cm}$.
    \end{enumerate}

    \begin{center}
        \adjustbox{max width=0.94\textwidth}{%
            \begin{tikzpicture}[>=Latex, scale=0.78]
                % SECCIÓN 1: Doble T
                \begin{scope}[shift={(0,0)}]
                    \node[above, font=\small\bfseries, blue!85!black] at (2.5, 5.8) {Sección 1 (Doble T)};
                    \fill[blue!12] (0, 0) rectangle (5.0, 1.0);
                    \fill[blue!12] (2.0, 1.0) rectangle (3.0, 4.0);
                    \fill[blue!12] (0, 4.0) rectangle (5.0, 5.0);
                    \draw[thick, blue!85!black] (0,0) -- (5,0) -- (5,1) -- (3,1) -- (3,4) -- (5,4) -- (5,5) -- (0,5) -- (0,4) -- (2,4) -- (2,1) -- (0,1) -- cycle;
                    \draw[->, dashed, red!80!black] (2.5, -0.4) -- (2.5, 5.6) node[right, font=\scriptsize\bfseries] {$y_G$};
                    \draw[->, dashed, red!80!black] (-0.5, 2.5) -- (5.05, 2.5);
                    \node[left=1pt, font=\scriptsize\bfseries, red!80!black] at (-0.5, 2.5) {$z_G$};
                    \draw[<->, line width=0.6pt] (0, -0.5) -- node[below=1pt, font=\tiny] {$20$} (2.0, -0.5);
                    \draw[<->, line width=0.6pt] (2.0, -0.5) -- node[below=1pt, font=\tiny] {$10$} (3.0, -0.5);
                    \draw[<->, line width=0.6pt] (3.0, -0.5) -- node[below=1pt, font=\tiny] {$20$} (5.0, -0.5);
                    \draw[<->, line width=0.6pt] (5.3, 0) -- node[right=1pt, font=\tiny] {$10$} (5.3, 1.0);
                    \draw[<->, line width=0.6pt] (5.3, 1.0) -- node[right=1pt, font=\tiny] {$30$} (5.3, 4.0);
                    \draw[<->, line width=0.6pt] (5.3, 4.0) -- node[right=1pt, font=\tiny] {$10$} (5.3, 5.0);
                \end{scope}

                % SECCIÓN 2: Cruz
                \begin{scope}[shift={(7.2,0)}]
                    \node[above, font=\small\bfseries, orange!85!black] at (2.5, 5.8) {Sección 2 (Cruz)};
                    \fill[orange!14] (2.0, 0) rectangle (3.0, 5.0);
                    \fill[orange!14] (0, 1.5) rectangle (5.0, 3.5);
                    \draw[thick, orange!85!black] (2,0) -- (3,0) -- (3,1.5) -- (5,1.5) -- (5,3.5) -- (3,3.5) -- (3,5) -- (2,5) -- (2,3.5) -- (0,3.5) -- (0,1.5) -- (2,1.5) -- cycle;
                    \draw[->, dashed, red!80!black] (2.5, -0.4) -- (2.5, 5.6) node[right, font=\scriptsize\bfseries] {$y_G$};
                    \draw[->, dashed, red!80!black] (-0.5, 2.5) -- (5.05, 2.5);
                    \node[left=1pt, font=\scriptsize\bfseries, red!80!black] at (-0.5, 2.5) {$z_G$};
                    \draw[<->, line width=0.6pt] (0, -0.5) -- node[below=1pt, font=\tiny] {$20$} (2.0, -0.5);
                    \draw[<->, line width=0.6pt] (2.0, -0.5) -- node[below=1pt, font=\tiny] {$10$} (3.0, -0.5);
                    \draw[<->, line width=0.6pt] (3.0, -0.5) -- node[below=1pt, font=\tiny] {$20$} (5.0, -0.5);
                    \draw[<->, line width=0.6pt] (5.3, 0) -- node[right=1pt, font=\tiny] {$15$} (5.3, 1.5);
                    \draw[<->, line width=0.6pt] (5.3, 1.5) -- node[right=1pt, font=\tiny] {$20$} (5.3, 3.5);
                    \draw[<->, line width=0.6pt] (5.3, 3.5) -- node[right=1pt, font=\tiny] {$15$} (5.3, 5.0);
                \end{scope}

                % SECCIÓN 3: Rectángulo
                \begin{scope}[shift={(14.4,0)}]
                    \node[above, font=\small\bfseries, green!55!black] at (1.3, 5.8) {Sección 3 (Rectángulo)};
                    \fill[green!12] (0, 0) rectangle (2.6, 5.0);
                    \draw[thick, green!60!black] (0, 0) rectangle (2.6, 5.0);
                    \draw[->, dashed, red!80!black] (1.3, -0.4) -- (1.3, 5.6) node[right, font=\scriptsize\bfseries] {$y_G$};
                    \draw[->, dashed, red!80!black] (-0.5, 2.5) -- (2.65, 2.5);
                    \node[left=1pt, font=\scriptsize\bfseries, red!80!black] at (-0.5, 2.5) {$z_G$};
                    \draw[<->, line width=0.6pt] (0, -0.5) -- node[below=1pt, font=\tiny] {$26\text{ cm}$} (2.6, -0.5);
                    \draw[<->, line width=0.6pt] (2.9, 0) -- node[right=1pt, font=\tiny] {$50\text{ cm}$} (2.9, 5.0);
                \end{scope}
            \end{tikzpicture}%
        }
    \end{center}

    \textbf{Determinar:}
    \begin{enumerate}[label=\alph*)]
        \item El máximo esfuerzo axil de tracción que puede resistir cada sección.
        \item El máximo momento flector positivo que puede resistir cada sección.
        \item La conclusión estructural que se extrae de la comparación.
    \end{enumerate}

    \tcblower

    \paso{1}{Apartado a) Área Transversal y Máximo Esfuerzo de Tracción}

    Comprobamos el área de las tres secciones transversales:
    \begin{align}
        \Omega_1 & = 2 \cdot (50 \cdot 10) + 10 \cdot 30 = 1000 + 300 = \mathbf{1300\text{ cm}^2} \\
        \Omega_2 & = 10 \cdot 50 + 2 \cdot (20 \cdot 20) = 500 + 800 = \mathbf{1300\text{ cm}^2}  \\
        \Omega_3 & = 26 \cdot 50 = \mathbf{1300\text{ cm}^2}
    \end{align}
    Las tres secciones poseen exactamente la \textbf{misma área transversal} $\Omega = 1300\text{ cm}^2 = 130\,000\text{ mm}^2$.
    Con la resistencia de cálculo del acero $f_{yd} = \frac{f_{yk}}{\gamma_s} = \frac{275}{1.05} = 261.905\text{ N/mm}^2$, el esfuerzo axil máximo de cálculo que agota la sección a tracción es idéntico para las tres:
    \begin{equation}
        \mathbf{N_d = \Omega \cdot f_{yd} = 130\,000\text{ mm}^2 \times 261.905\text{ N/mm}^2 = 34\,047\,619\text{ N} \approx 34\,047.6\text{ kN}}
    \end{equation}
    (En valor característico de servicio: $N_k = \frac{N_d}{\gamma_f} = \frac{34\,047.6}{1.5} = 22\,698.4\text{ kN}$).

    \paso{2}{Apartado b) Momentos de Inercia $I_{z_G}$ y Momento Flector Máximo $M_d$}

    Al tener todas las secciones doble simetría y canto total $H = 50\text{ cm}$, la distancia desde la fibra neutra hasta la fibra extrema es común e igual a $y_{\text{máx}} = 25\text{ cm} = 250\text{ mm}$. Calculamos el momento de inercia $I_{z_G}$ y el momento último $M_d = \frac{f_{yd} \cdot I_{z_G}}{y_{\text{máx}}}$ para cada sección:
    \begin{enumerate}
        \item \textbf{Sección 1 (Doble T):} Calculado por diferencia entre el rectángulo envolvente ($50 \times 50$) y los dos huecos laterales ($2 \times 20 \times 30$):
              \begin{equation}
                  I_{z_G,1} = \frac{50 \cdot 50^3}{12} - 2 \cdot \frac{20 \cdot 30^3}{12} = 520\,833.33 - 90\,000 = \mathbf{430\,833.33\text{ cm}^4}
              \end{equation}
              \begin{equation}
                  \mathbf{M_{d,1} = \frac{261.905\text{ N/mm}^2 \cdot \left(430\,833.33 \times 10^4\text{ mm}^4\right)}{250\text{ mm}} = 4513.49 \times 10^6\text{ N}\cdot\text{mm} \approx 4513.4\text{ kN}\cdot\text{m}}
              \end{equation}

        \item \textbf{Sección 2 (Cruz):} Sumando la inercia del nervio vertical ($10 \times 50$) y de las dos alas centrales ($40 \times 20$):
              \begin{equation}
                  I_{z_G,2} = \frac{10 \cdot 50^3}{12} + 2 \cdot \frac{20 \cdot 20^3}{12} = 104\,166.67 + 26\,666.67 = \mathbf{130\,833.33\text{ cm}^4}
              \end{equation}
              \begin{equation}
                  \mathbf{M_{d,2} = \frac{261.905\text{ N/mm}^2 \cdot \left(130\,833.33 \times 10^4\text{ mm}^4\right)}{250\text{ mm}} = 1370.63 \times 10^6\text{ N}\cdot\text{mm} \approx 1370.6\text{ kN}\cdot\text{m}}
              \end{equation}

        \item \textbf{Sección 3 (Rectángulo $26 \times 50$):}
              \begin{equation}
                  I_{z_G,3} = \frac{26 \cdot 50^3}{12} = \mathbf{270\,833.33\text{ cm}^4}
              \end{equation}
              \begin{equation}
                  \mathbf{M_{d,3} = \frac{261.905\text{ N/mm}^2 \cdot \left(270\,833.33 \times 10^4\text{ mm}^4\right)}{250\text{ mm}} = 2837.30 \times 10^6\text{ N}\cdot\text{mm} \approx 2837.2\text{ kN}\cdot\text{m}}
              \end{equation}
    \end{enumerate}

    \paso{3}{Apartado c) Conclusión de Diseño Estructural}

    \begin{tcolorbox}[colback=yellow!10, colframe=orange!85!black, arc=3pt, boxrule=0.8pt]
        \textbf{Conclusión Fundamental:}
        \begin{itemize}
            \item \textbf{Frente a Esfuerzo Axil ($N$):} Las 3 secciones tienen exactamente la \textbf{misma área ($\Omega = 1300\text{ cm}^2$)} y, por tanto, consumen la misma cantidad de acero y \textbf{aguantan exactamente el mismo esfuerzo axil ($34\,047.6\text{ kN}$)}.
            \item \textbf{Frente a Momento Flector ($M$):} Poseen \textbf{muy distinto momento de inercia ($I_z$)}. A mayor momento de inercia, mayor momento flector resiste la pieza: el perfil en \textbf{Doble T} ($4513.4\text{ kN}\cdot\text{m}$) resiste un $\mathbf{59\%}$ más que el rectángulo macizo ($2837.2\text{ kN}\cdot\text{m}$) y $\mathbf{3.29\text{ veces}}$ más que la sección en cruz ($1370.6\text{ kN}\cdot\text{m}$) empleando el mismo peso de material.
        \end{itemize}
    \end{tcolorbox}
}
